\documentclass[11pt,letterpaper]{article}
\usepackage[margin=0.85in,headheight=14pt]{geometry}
\usepackage{amsmath,amssymb}
\usepackage{enumitem}
\usepackage{fancyhdr}
\usepackage{lmodern}
\usepackage[T1]{fontenc}
\usepackage[dvipsnames]{xcolor}
\usepackage[most]{tcolorbox}

\pagestyle{fancy}
\fancyhf{}
\lhead{SYDE 556/750 \qquad Recipe practice}
\rhead{Feedforward transformations}
\cfoot{\thepage}
\renewcommand{\headrulewidth}{0pt}

\setlength{\parindent}{0pt}
\setlength{\parskip}{4pt}

\newcommand{\Jb}{J^{\mathrm{bias}}}

\newtcolorbox{idea}{colback=gray!8,colframe=gray!55,boxrule=0.4pt,left=5pt,right=5pt,top=2pt,bottom=2pt,arc=1pt,fonttitle=\bfseries\small,title=The recipe}

\setlist[enumerate]{leftmargin=1.6em,itemsep=6pt,topsep=4pt}

\begin{document}

\begin{center}
{\Large\bfseries Practice set: the feedforward recipe}\\[2pt]
{\large Lecture 5, one transformation at a time}\\[4pt]
{\normalsize Questions. Solutions are in a separate file. Leave that file closed.}
\end{center}

\subsection*{How to use this}

Each question is the four-step recipe from the lecture slide. Do the steps in order. Question 2 keeps the populations and changes only the function, so a memorized weight matrix from question 1 will not survive.

\begin{idea}
\begin{enumerate}[leftmargin=1.4em,itemsep=3pt]
\item Write the encoding of the input population and the encoding of the output population.
\item Write the transformation in the variables those populations represent.
\item Solve for the function decoder $D^f$ of the input population. With no noise penalty,
\[
(D^f)^{\mathsf T} = (AA^{\mathsf T})^{-1} A Y^{\mathsf T},
\]
where column $k$ of $Y$ is $f(x_k)$. The noise-aware formula is the same replacement of $X$ by $Y$.
\item Substitute $\hat y = D^f a$ into the output encoding. With row $i$ of $E$ equal to $\alpha_i e_i^{\mathsf T}$,
\[
W^f = E D^f, \qquad J = W^f a + \Jb.
\]
The bias is added once. It is not a column of $W^f$.
\end{enumerate}
\end{idea}

$A$ has one row per input neuron and one column per training sample. Leave fractions as fractions.

%---------------------------------------------------------------------
\section{One pass through all four steps}

Population A represents a scalar $x$ with two neurons. At the training inputs $x = (0,\ 1)$ their activities are
\[
a_1 = (1,\ 2), \qquad a_2 = (2,\ 1).
\]
Neuron 1 of A has $\alpha = 1$, $e = +1$, $\Jb = 1$. Neuron 2 of A has $\alpha = 1$, $e = -1$, $\Jb = 1$.

Population B represents a scalar $y$ with two neurons.
Neuron 1 of B has $\alpha = 2$, $e = +1$, $\Jb = 1/2$.
Neuron 2 of B has $\alpha = 1$, $e = -1$, $\Jb = 1$.

The connection computes $y = 2x$. Use the unregularized decoder.

\begin{enumerate}
\item Write $J$ for each neuron of A in terms of $x$, and $J$ for each neuron of B in terms of $y$.
\item Write the transformation the connection is supposed to compute.
\item Form $A$ and $Y$. Compute $D^f$. Check $D^f a$ against $2x$ on both training inputs.
\item Substitute the decoder into B's encoding. Give $W^f$, then write both currents as weighted sums of $a_1$ and $a_2$. At $x = 1$, compute both currents from those weights and again from $\alpha\, e\, (2x) + \Jb$. The two routes should agree.
\end{enumerate}

%---------------------------------------------------------------------
\section{Same populations, new function}

The populations, activities, and training inputs stay as in the previous question. The connection now computes $y = x + 1$.

\begin{enumerate}
\item Say which of the four steps can be copied from the previous question, and which must be redone.
\item Compute the new $D^f$. The inverse of $AA^{\mathsf T}$ can be reused.
\item Substitute again. Write both currents of B.
\item An identity decoder for this $A$ is $d = (2/3,\ -1/3)$. A decoder for the constant function $1$ is $d^{\mathbf 1} = (1/3,\ 1/3)$. Show that $d + d^{\mathbf 1}$ equals the decoder you just found, and say what $y$ the identity decoder alone would implement.
\end{enumerate}

\end{document}
